% C.3bc. Final proof, independently audited by root and real_weil.
\paragraph{C.3bc. Top Chow groups for SNC resolutions over
\(\overline{\mathbb F}_p\) (Proved).}
Let \(k=\overline{\mathbb F}_p\), and let \(R\) be the strict
henselization of a normal local ring essentially of finite type over
\(k\), at a closed point with residue field \(k\).
Write \(\dim R=d+1\), where \(d\geq1\).
Suppose that \(f:Z\to\operatorname{Spec}R\) is projective and
birational, that \(Z\) is regular, and that \(f\) is an isomorphism
away from the closed point. Suppose that its reduced exceptional
fibre is a simple normal crossing divisor
\[
                       E=\bigcup_{i=1}^{c}E_i.
\]
Assume that there are positive rational numbers \(a_i\) such that
the exceptional \(\mathbb Q\)-Cartier divisor
\[
                    D=\sum_i a_iE_i
\]
satisfies \(-D|_E\) ample. Put \(X=Z\setminus E\), and let
\(i_i:E_i\hookrightarrow Z\) denote the inclusions. Then
\begin{equation}\label{eq:C3bc-component-degrees}
 \mathrm{CH}^{d}(Z)_{\mathbb Q}
   \xrightarrow[\alpha\mapsto(\deg i_i^!\alpha)_i]{\ \sim\ }
                 \mathbb Q^{c},
 \qquad
                 \mathrm{CH}^{d}(X)_{\mathbb Q}=0.
\end{equation}
In particular the ordinary quotient \(B^d(X)_{\mathbb Q}\) is zero.
The statement assumes the indicated resolution and ample exceptional
combination; it does not assume their existence in arbitrary dimension.
No lifting of algebraic families from the reducible exceptional fibre
is used.

\begin{proof}
\emph{1. The top weight on a smooth stratum.}
We fix an integer \(q>1\), and use the generalized \(q^w\)-eigenspace
of the Adams operation \(\psi^q\) on rational \(K\)-groups as the
weight-\(w\) part, denoted by a superscript \((w)\).
For a smooth quasi-projective \(k\)-scheme \(V\) of pure dimension
\(m\), the higher Chow comparison gives
\[
                  K_n(V)_{\mathbb Q}^{(w)}
                    \simeq\mathrm{CH}^{w}(V,n)_{\mathbb Q}
                         \qquad(n\geq0).
\]
We use \citet[Theorem 3.1, pp.~276--278]{frC3bcLevine1994};
its cubical and simplicial comparisons are equivalent by that
paper's Theorem 4.7.
Since \(V\) times the \(n\)-cube has dimension \(m+n\), the right
side vanishes for \(w>m+n\).
For \(w=m+n\), its degree-\(n\) cycles are zero-cycles in the
interior of \(V\) times the cube. Every point is \(k\)-rational.
Projecting it to \(x\in V(k)\) shows that it is the proper
pushforward of a degree-\(n\) cycle on the cube over
\(\operatorname{Spec}k\). The degree-\((n-1)\) cycle group is
zero, so all these generators are cycles. Proper pushforward
respects their boundaries. Consequently
\[
 \bigoplus_{x\in V(k)}\mathrm{CH}^{n}(k,n)
        \longrightarrow\mathrm{CH}^{m+n}(V,n)
                   \quad\hbox{is surjective}.
\]
By \citet[Theorem 1, p.~181]{frC3bcTotaro1992},
\(\mathrm{CH}^{n}(k,n)=K_n^M(k)\).
For \(n>0\) this is torsion: the Milnor group is a quotient of
the \(n\)-fold tensor power of \(k^*\), and every element of
\(k^*\) has finite order. We have proved
\begin{equation}\label{eq:C3bc-smooth-weights}
 \begin{aligned}
 K_n(V)_{\mathbb Q}^{(w)}&=0
                 &&(w>m+n,\ n\geq0),\\
 K_n(V)_{\mathbb Q}^{(m+n)}&=0
                 &&(n>0).
 \end{aligned}
\end{equation}
If \(V\) is smooth projective and connected, then
\(\mathrm{CH}_0(V)_{\mathbb Q}=\mathbb Q\) by degree.
Indeed, two points lie on a smooth connected projective curve;
the prescribed-point Bertini argument works in every characteristic
\citep[Lemma 45.8.3, proof]{frC3azStacksBertini}.
Their difference in the curve's Jacobian is defined over a finite
field and is therefore torsion. Pushing it forward proves the
assertion. The zero-dimensional case is immediate.

\emph{2. Mayer--Vietoris without an Adams weight shift.}
Let \(V=A\cup B\) be a union of components of a projective SNC
scheme, and put \(C=A\cap B\) with its reduced scheme structure.
Locally the coordinate ring is a fibre product of the rings of
\(A\) and \(B\) over that of \(C\), with surjective arrows.
The bi-relative excision theorem of
\citet[Theorem C, p.~30]{rootGeisserHesselholt2011}
therefore makes the ordinary \(K\)-theory square homotopy
cartesian after inverting \(p\).
Finite affine Mayer--Vietoris descent
\citep[Theorem 8.1, p.~367]{D_ThomasonTrobaugh1990}
patches this to the projective schemes. In particular we have
the rational exact sequence
\begin{equation}\label{eq:C3bc-MV}
 K_{n+1}(C)_{\mathbb Q}\longrightarrow K_n(V)_{\mathbb Q}
 \longrightarrow K_n(A)_{\mathbb Q}\oplus K_n(B)_{\mathbb Q}
 \longrightarrow K_n(C)_{\mathbb Q}.
\end{equation}
Equivalently, one may apply KST to the closed-component abstract
blow-up square and use nilpotent invariance after inverting \(p\).

All arrows in the square are ordinary pullbacks of vector bundles.
The connecting arrow is the boundary of their iterated homotopy
fibre. Here we use the paragraph \emph{before} Lemma 2.1 on
p.~258 of \citet{frC3bcLevine1994}: it describes Grayson's
exact-category construction for arbitrary schemes, its extension
to any iterated homotopy fibre, and its delooping. Pullback is exact
on vector bundles and commutes with their tensor, symmetric, and
exterior powers. This construction therefore gives a natural
endomorphism of the closed-union square and its fibre sequence,
including the boundary. It agrees with the usual operations on
the smooth strata and on \(K_0\). Thus \eqref{eq:C3bc-MV}
preserves \(w\).
We do not identify a group with supports with a group on a closed
subscheme by d\'evissage, so there is no codimension or Thom shift
in this assertion. The finite generalized-eigenspace decomposition
is preserved under subquotients and extensions. We use only the
elementary \(\mathbb Q[\psi^q]\)-module argument in the
\emph{proof} of Lemma 2.1, p.~259, at this point, not its
regular-ambient statement for the singular scheme \(V\).
It follows inductively from the smooth strata that taking a fixed
generalized weight in \eqref{eq:C3bc-MV} is exact.

\emph{3. The top weight of an SNC scheme.}
Induction on dimension and on the number of components proves,
for a projective SNC scheme \(V\) of pure dimension \(m\),
\begin{equation}\label{eq:C3bc-SNC-weights}
 \begin{aligned}
 K_n(V)_{\mathbb Q}^{(w)}&=0
                 &&(w>m+n,\ n\geq0),\\
 K_n(V)_{\mathbb Q}^{(m+n)}&=0
                 &&(n>0),\\
 K_0(V)_{\mathbb Q}^{(m)}
       &\xrightarrow{\ \sim\ }
           \bigoplus_{V_j}\mathrm{CH}_0(V_j)_{\mathbb Q}.
 \end{aligned}
\end{equation}
Here \(V_j\) are its irreducible components, and the last map is
restriction followed by the rational Chern character.
For completeness, in the induction step take \(B\) to be one
component. Then \(C\) is SNC of dimension \(m-1\), with possibly
several connected components. For \(n>0\), the predecessor of
weight \(m+n\) in \eqref{eq:C3bc-MV} is
\[
                  K_{n+1}(C)_{\mathbb Q}^{(m+n)},
\]
which is exactly the top weight
\((m-1)+(n+1)\) of a positive \(K\)-group and vanishes by the
dimension induction. We obtain the second line; the same argument
above that weight gives the first line. For \(n=0,w=m\), the two
end terms are \(K_1(C)_{\mathbb Q}^{(m)}=0\) and
\(K_0(C)_{\mathbb Q}^{(m)}=0\), respectively. This gives the last
line by the component induction. The base case consists of smooth
strata and \eqref{eq:C3bc-smooth-weights}. Empty intersections
cause no change.
In particular,
\begin{equation}\label{eq:C3bc-E-degrees}
             K_0(E)_{\mathbb Q}^{(d)}
                      \xrightarrow{\ \sim\ }\mathbb Q^c.
\end{equation}

\emph{4. Finite detection on the resolution.}
The finite-detection argument of C.3ax and C.3az applies without
smoothness or irreducibility of \(E\). To recall it, set
\(E_b=Z\times_RR/\mathfrak m_R^b\). KST
\citep[Theorem A and (1.6), author pp.~2--3]
{rootKerzStrunkTamme2018} gives the exact pro sequence
\[
 K_0(R)\longrightarrow
 \{K_0(R/\mathfrak m_R^b)\}_b\oplus K_0(Z)
             \longrightarrow\{K_0(E_b)\}_b.
\]
The two local-ring groups are \(\mathbb Z\); quotient by the
graph of the rank-one structure sheaf. The resulting map from
the constant pro group \(K_0(Z)\) is monic. Its kernels have
inclusion transition maps, so pro-zero implies that one finite
kernel is zero. Thus \(K_0(Z)\to K_0(E_{b_0})\) is injective
for some \(b_0\). Since \((E_b)_{\mathrm{red}}=E\), GH
\citep[Theorem A, p.~29]{rootGeisserHesselholt2011} and finite
affine descent give
\[
                         K_0(Z)_{\mathbb Q}
                              \hookrightarrow K_0(E)_{\mathbb Q}.
\]
Descend the finite data to smooth pointed-\'etale models of \(Z\)
as in C.3az. The rational Chern ring isomorphism
\citep[Example 15.2.16(b), pp.~294--295]{frC4Fulton1998}
then identifies \(\mathrm{CH}^d(Z)_{\mathbb Q}\) with its
weight-\(d\) part, compatibly with restriction to each smooth
\(E_i\). By \eqref{eq:C3bc-E-degrees} this proves injectivity
of the first arrow in \eqref{eq:C3bc-component-degrees}.
This step does not require a comparison between \(K\) and \(KH\).

\emph{5. Exceptional curves span all component degrees.}
Choose an ample line bundle \(H\) on \(E\), and put
\[
 C_j=c_1(H|_{E_j})^{d-1}\cap[E_j]\in\mathrm{CH}_1(E_j).
\]
For \(d=1\) this means \(C_j=[E_j]\).
Regard each \(C_j\) as an exceptional boundary class in
\(\mathrm{CH}^d(Z)\). Its component-degree vector is the \(j\)-th
column of the rational symmetric matrix
\[
 M_{ij}
   =\deg\!\left(c_1(\mathcal O_Z(E_i)|_{E_j})
                       c_1(H|_{E_j})^{d-1}\right).
\]
For \(i\ne j\), this is the degree of \(H^{d-1}\) on
\(E_i\cap E_j\), and is nonnegative. The common line bundle \(H\)
on \(E\) also proves symmetry. If \(a=(a_i)_i\), then
\[
 (Ma)_i
      =\deg\!\left(c_1(D|_{E_i})c_1(H|_{E_i})^{d-1}\right)<0
\]
because \(-D|_{E_i}\) is ample.
Set \(N=\operatorname{diag}(a)M\operatorname{diag}(a)\) and
\(b_i=\sum_jN_{ij}=a_i(Ma)_i<0\).
For every \(x\in\mathbb R^c\) we have the exact identity
\[
 x^{\mathsf t}Nx
    =\sum_i b_ix_i^2
       -\sum_{i<j}N_{ij}(x_i-x_j)^2.
\]
Hence \(N\), and therefore \(M\), is negative definite and
invertible. The exceptional classes \(C_j\) span \(\mathbb Q^c\)
under the injective degree map, proving that map is an
isomorphism. Finally Chow localization is surjective from
\(\mathrm{CH}^d(Z)_{\mathbb Q}\) to
\(\mathrm{CH}^d(X)_{\mathbb Q}\), and every spanning class
\(C_j\) is supported on \(E\). The latter group is zero.
\end{proof}
